\honest{My Childhood Role Model}{Eliezer Yudkowsky}{2008}

When I lecture on the Singularity, I often draw a scale of intelligence with the ``village idiot'' at one end and Einstein at the other. That is the parochial view of everyday life, where we only meet human minds; for AI we need the scale of all possible minds. On the scale of possible minds, the two are almost the same point: ``Einstein and the village idiot both have a prefrontal cortex, a hippocampus, a cerebellum\ldots'' A chimpanzee is much farther away. A chimp could not tell Einstein from the village idiot, and our descendants may not either. Carl Shulman has noticed that some academics who talk about transhumanism put Einstein far above the village idiot. \nb{Brain size, which the post does not cite, fits its ordering.}

At the 2006 Singularity Summit, Douglas Hofstadter looked at my diagram and said, ``I think Einstein should be way off on the right.'' I was speechless, all the more since Hofstadter was one of my childhood heroes, and took it as a ``cultural gap,'' since I met the idea of a Jupiter Brain at twelve. I do not give Hofstadter's reasons. The right of my scale holds Jupiter Brains; Douglas Adams's Deep Thought, which, switched on for the first time, started from ``I think therefore I am'' and deduced rice pudding and income tax before anyone could shut it off; the Elders of Arisia, galactic overminds, Matrioshka brains and the better class of God; at the far end, the Old One and the Blight. A General Systems Vehicle would find Einstein very cute. ``Not frickin' Einstein.''

I admit this is generalization from fictional evidence. Still, I suspect science fiction helps people imagine outside their own era, instead of a world where humans always have existed and always will. \nb{Yudkowsky's ``The Logical Fallacy of Generalization from Fictional Evidence,'' earlier in the book, said the most damaging thing about other authors' imaginations ``is that it stops people from using their own.''} The universe is 13.7 billion years old, and our species ``a hundred thousand years or thereabouts.'' \nb{Fossils dated in 2005 and 2017 put it at about 195,000 and 315,000 years. The point survives.} Then again, some who never read science fiction can imagine other worlds, and some fans cannot.

Yesterday I wanted to say that Einstein, cute for a human, was about as efficient with evidence as the US Department of Defense. I depicted a civilization of Einsteins because a superintelligence with a webcam would be imagined as a ghost in a box that has not been told how to interpret pictures, so does not know; readers would not apply their own creativity to the problem. Could one superintelligence do it all? Consider that Jeffreyssai counted 5,760 five-minute insights in a month. Human brains were adapted to chase deer, spear and cook them, and argue cleverly for a larger share of the meat, so repurposing one for physics deserves a record, like the fastest car ever built of Jello. The full horror of how the blind idiot god designed the brain dawns only with much cognitive science, and the biases are a hint. We need full concentration to multiply three-digit numbers. Einstein used his neurons as inefficiently as his data. \nb{In ``The Failures of Eld Science,'' Jeffreyssai himself says that rate is not realistic. The list shows that human brains are inefficient, not how much better one mind could do.}

I have ``certain ulterior motives,'' which I do not name. My childhood ideal was never Einstein but a dream of an AI. Ideals, Carl Schurz said, are like stars. The dream ``helped me to some degree, and harmed me to some degree,'' and I do not say how. \nb{In ``My Childhood Death Spiral,'' later in the book, Yudkowsky describes an early admiration of intelligence that led to ``false happy beliefs,'' above all that a superintelligence would be moral. That post does not say this is the harm meant here.} Some ideals are like dreams, from within us, as Mentor of Arisia came from E. E. Smith's imagination; I grant that an imagined ideal ``is only your own mind talking,'' and if you guess wrong you go astray.

But do not take your ideals only from real, dead humans. Each generation can do better, and asking ``Do I dare to do this thing, which Einstein could not do?'' is like Einstein asking whether he may do better than Newton. Following stars, ``at best, it gets you to the star.'' Your era supports you more than you know, in assumptions and technology of mind. \nb{Einstein revered Newton and went past him: ``Newton, forgive me,'' he wrote in 1949.} Einstein ``talked a deal of nonsense about an impersonal God.'' I quote none of it. \nb{Einstein rejected a personal God and used the word for the order of nature: ``Spinoza's God who reveals himself in the orderly harmony of what exists.''}

It seems less like sacrilege with a galactic supermind on your scale beside Einstein. If you try only what seems humanly possible, you ask too little of yourself, and reasons why the higher goal is ``not possible'' leap to mind. The most important role models are dreams of perfection; dreaming of less draws on less than the full power of the part of you that dreams.
